\Subhead{2.6.2}{Диэлектрические волноводы. Световоды}

В диэлектрическом волноводе поле удерживается большим показателем сердцевины, а в оболочке затухает. Для плоской сердцевины $n_1>n_2$ направляемая мода имеет
\begin{equation}
 n_2k_0<\beta<n_1k_0,\quad
 h^2=n_1^2k_0^2-\beta^2>0,\quad
 \alpha^2=\beta^2-n_2^2k_0^2>0.
\end{equation}
В сердцевине поперечное поле осциллирует, в оболочке экспоненциально затухает. Граничная непрерывность касательных $E,H$ квантует $\beta$. Например, для симметричной TE-моды плоского слоя толщины $d$ имеем $h\tan(hd/2)=\alpha$. Это волновая версия полного внутреннего отражения; лучевое объяснение не определяет точные дисперсии и отсечки.

Для ступенчатого круглого волокна радиуса $a$ параметр
\begin{equation}
 V=k_0a\sqrt{n_1^2-n_2^2}
\end{equation}
измеряет число поперечных осцилляций. При слабой направляющей $n_1-n_2\ll n_1$ фундаментальная HE$_{11}$-мода не имеет ненулевой отсечки; условие $V<2.405$ обеспечивает отсутствие следующего семейства направляемых мод в идеальной ступенчатой модели. Числовая апертура в воздухе приблизительно $\sqrt{n_1^2-n_2^2}$. Потери, изгибы и межмодовая дисперсия ограничивают передачу; материал и геометрия оба вносят вклад в дисперсию одной моды.
